\section{Finite Spatial Correlation Kernels}
\label{sec:ch19-spatial-correlators}

A density mode is a useful algebraic coordinate, but a physical density is
usually discussed at a \emph{site}.  To pass from the ordered mode
contractions of Section~\ref{sec:ch19-mode-correlators} to an observable
at two positions, we must reconstruct the spatial field by a finite Fourier
sum.  This operation appears innocuous: multiply each mode by a phase and
sum.  Yet the order of the two resulting fields matters, because each field
contains both creation and annihilation generators.  Reversing their product
changes the answer by an exact canonical commutator.

The question of this section is therefore deliberately finite and precise:
\emph{given the positive abstract CCR functional constructed in
Section~\ref{sec:ch19-scope}, what complex number does it assign to the
ordered product of two spatial density operators?}  We will first derive
that number, then separate its symmetric (fluctuation) and antisymmetric
(commutator) contributions.  Only afterward will we interpret positivity
and evaluate a concrete lattice example.  Neither this reconstruction nor
the finite Fourier calculation identifies the abstract functional with a
ground-state expectation on a compressed fermionic energy-budget space.

\subsection{A finite lattice and its Fourier phases}

In the source implementation, the lattice length is \emph{even}:
\begin{equation}
 L=2h,\qquad h\in\mathbb N,\quad h>0,\qquad
 \Lambda_L=\mathbb Z/L\mathbb Z.
 \label{eq:ch19-spatial-lattice}
\end{equation}
A point $x\in\Lambda_L$ is a residue class, represented when necessary
by its unique integer $x_0\in\{0,\ldots,L-1\}$.  Choose the canonical
primitive $L$th root
\begin{equation}
 \zeta=\exp(2\pi i/L),\qquad \zeta^L=1,
 \quad |\zeta|=1.
 \label{eq:ch19-spatial-zeta}
\end{equation}
More generally, every result written in terms of residue-class
characters holds when $\zeta$ is \emph{any} primitive $L$th root.
The word \emph{primitive} means that $\zeta^k\ne1$ for
$1\leq k<L$: the powers of $\zeta$ distinguish the $L$ sites.
The source also proves some intermediate algebraic identities for a
merely nonzero complex number $\zeta$.  That weaker hypothesis suffices
for powers involving the chosen representatives $x_0,y_0$, but the
periodic character interpretation and the Hermiticity arguments below
use the root-of-unity condition.

For $m\in\mathbb Z$ define the character
\begin{equation}
 \chi_m(x):=\zeta^{m x_0}.
 \label{eq:ch19-spatial-character}
\end{equation}
This is well defined on residue classes because replacing $x_0$ by
$x_0+kL$ leaves its value unchanged.  Moreover,
\begin{equation}
 \chi_m(x+y)=\chi_m(x)\chi_m(y),\qquad
 \chi_{-m}(x)=\overline{\chi_m(x)}=\chi_m(x)^{-1},\qquad
 \chi_m(x-y)=\chi_m(x)\chi_{-m}(y).
 \label{eq:ch19-spatial-character-laws}
\end{equation}
The last identity is the step that turns a double Fourier sum into a
function of the \emph{separation} $x-y$, rather than of the two
positions independently.  Notice that the separation belongs to
$\Lambda_L$; no choice of an infinite-lattice coordinate is needed.

Fix a positive-mode cutoff $M\in\mathbb N$ and retain the indices
$m\in\{1,\ldots,M\}$.  The source algebraic theorem allows $M=0$
(the sum is then empty), and does not assume $M<L/2$.  If one wants
these labels to represent distinct nonzero lattice frequencies away
from the Nyquist mode, the extra condition $M<h$ is natural, but it is
\emph{not a hypothesis} of the abstract kernel identity.  For $M\geq L$
some characters repeat even though the abstract oscillator generators
remain distinct.  The number of modes $M$ is also not an occupation
cutoff: each mode still acts on the uncompressed polynomial carrier.

\subsection{The spatial density as a Hermitian algebra element}

Let $C_{Rm}$ and $A_{Rm}=C_{Rm}^{*}$ denote the right-branch bare
oscillator generators, with exact uncompressed CCR
\begin{equation}
 [A_{Rm},C_{Rn}]=m\delta_{mn}I,\qquad
 [A_{Rm},A_{Rn}]=[C_{Rm},C_{Rn}]=0.
 \label{eq:ch19-spatial-ccr}
\end{equation}
We identify the positive- and negative-index density modes as
$\rho_{m,R}=C_{Rm}$ and $\rho_{-m,R}=A_{Rm}$, respectively.
Define the right-branch \emph{oscillatory density} by
\begin{equation}
 \rho_R(x):=\frac1L\sum_{m=1}^{M}
 \bigl(\chi_{-m}(x)C_{Rm}+\chi_m(x)A_{Rm}\bigr).
 \label{eq:ch19-spatial-rho}
\end{equation}
It contains no charge zero mode.  This convention uses $1/L$ rather
than $1/\sqrt L$; the normalizations of the Fourier transform and of
the oscillator weights should not be silently combined.  In
particular, the cutoff $M$ labels a selected collection of oscillatory
modes and does not turn \eqref{eq:ch19-spatial-rho} into a claim about
the full microscopic lattice density.

Why is this definition an observable rather than an arbitrary complex
combination?  The conjugate character multiplies the adjoint operator:
\begin{align*}
 \rho_R(x)^*
 &=\frac1L\sum_{m=1}^{M}
   \bigl(\overline{\chi_{-m}(x)}A_{Rm}
        +\overline{\chi_m(x)}C_{Rm}\bigr)\\
 &=\frac1L\sum_{m=1}^{M}
   \bigl(\chi_m(x)A_{Rm}+\chi_{-m}(x)C_{Rm}\bigr)
 =\rho_R(x).
\end{align*}
The reality of $1/L$ and the unit modulus of $\zeta$ were used here.
This simple check will be essential when we deduce positive
semidefiniteness of the correlation matrix; a generic nonzero
complex $\zeta$ would not justify the same argument.

Let $P=(c,s)$ be a real hyperbolic pair,
\begin{equation}
 c,s\in\mathbb R,\qquad c^2-s^2=1.
 \label{eq:ch19-spatial-hyperbolic}
\end{equation}
We evaluate in the normalized positive \emph{abstract} functional
$\omega_P=\omega_0\circ\beta_{c,s}^{-1}$ of the preceding sections.
Its mode contractions, in the exact source ordering, are
\begin{align}
 \omega_P(C_{Rm}A_{Rn})&=m\delta_{mn}s^2,
 &\omega_P(A_{Rm}C_{Rn})&=m\delta_{mn}c^2,
 \label{eq:ch19-spatial-input-contractions}\\
 \omega_P(C_{Rm}C_{Rn})&=0,
 &\omega_P(A_{Rm}A_{Rn})&=0.
 \label{eq:ch19-spatial-input-zeros}
\end{align}
These are recalled as \emph{inputs}; their Bogoliubov derivation is not
repeated here.  The notation $\omega_P$ is deliberately distinct from
a finite-Fock vector expectation $\langle\Psi,\cdot\,\Psi\rangle$
or a density-matrix expectation $\operatorname{Tr}(\varrho\,\cdot)$.

\subsection{The exact ordered two-point kernel}

\begin{theorem}[Ordered spatial density kernel]
\label{thm:ch19-spatial-ordered}
Assume \eqref{eq:ch19-spatial-lattice}--\eqref{eq:ch19-spatial-hyperbolic}
and a nonnegative integer cutoff $M$.  For every $x,y\in\Lambda_L$,
the ordered correlation
$D_P(x,y):=\omega_P\bigl(\rho_R(x)\rho_R(y)\bigr)$ equals
\begin{equation}
 \boxed{\displaystyle
 D_P(x,y)=\frac1{L^2}\sum_{m=1}^{M}m
 \bigl(s^2\chi_{-m}(x-y)+c^2\chi_m(x-y)\bigr).}
 \label{eq:ch19-spatial-ordered}
\end{equation}
For the canonical root $\zeta=e^{2\pi i/L}$, the characters
in this formula are the exponentials
$\chi_{\pm m}(x-y)=e^{\pm2\pi i m(x_0-y_0)/L}$.
The value is independent of the chosen representatives.
\end{theorem}

\begin{proof}
Write $u_m(x)=\chi_{-m}(x)$ and $v_m(x)=\chi_m(x)$.
Expanding the product of the two finite sums \emph{before}
applying the functional gives
\begin{align*}
 L^2\rho_R(x)\rho_R(y)
 &=\sum_{m=1}^{M}\sum_{n=1}^{M}
  \bigl[u_m(x)u_n(y)C_{Rm}C_{Rn}\\
 &\qquad+u_m(x)v_n(y)C_{Rm}A_{Rn}\\
 &\qquad+v_m(x)u_n(y)A_{Rm}C_{Rn}\\
 &\qquad+v_m(x)v_n(y)A_{Rm}A_{Rn}\bigr].
\end{align*}
Because the sum is finite and $\omega_P$ is complex linear,
we may evaluate its four kinds of terms individually.  The $CC$
and $AA$ contributions vanish by
\eqref{eq:ch19-spatial-input-zeros}; the remaining two retain their
original order:
\begin{align*}
 L^2D_P(x,y)
 &=\sum_{m,n=1}^{M}u_m(x)v_n(y)
       \underbrace{\omega_P(C_{Rm}A_{Rn})}_{m\delta_{mn}s^2}\\
 &\qquad+\sum_{m,n=1}^{M}v_m(x)u_n(y)
       \underbrace{\omega_P(A_{Rm}C_{Rn})}_{m\delta_{mn}c^2}\\
 &=\sum_{m=1}^{M}m
  \bigl(s^2\chi_{-m}(x)\chi_m(y)
       +c^2\chi_m(x)\chi_{-m}(y)\bigr).
\end{align*}
The Kronecker delta removed one summation index.  Finally, the
character multiplication law
\eqref{eq:ch19-spatial-character-laws} replaces each product of two
phases by the corresponding character of $x-y$.  Dividing by $L^2$
proves the claim.  No infinite sum, reordering limit, or continuum
approximation enters the calculation.
\end{proof}

The notation $D_P(x,y)$ emphasizes the parameter dependence.  The
original chapter notes also call this quantity $D_c(x,y)$; that
subscript should not be misread as dependence on $c$ alone, since
both $c^2$ and $s^2$ occur.  The \emph{ordered} kernel is generally
complex for $x\neq y$.  Replacing the two exponentials immediately by
a cosine would discard its antisymmetric part and calculate a
different observable.

\subsection{The state-independent commutator and the symmetrized kernel}

For two algebra elements, write
$[B,C]=BC-CB$ and $\{B,C\}=BC+CB$.  By linearity,
\begin{equation}
 D_P(x,y)=\frac12\omega_P(\{\rho_R(x),\rho_R(y)\})
          +\frac12\omega_P([\rho_R(x),\rho_R(y)]).
 \label{eq:ch19-spatial-decomposition}
\end{equation}
The first term describes symmetric two-point fluctuations, while the
second remembers the operator ordering.  They are not interchangeable.

\begin{proposition}[Exact density commutator]
\label{prop:ch19-spatial-commutator}
Under the hypotheses of Theorem~\ref{thm:ch19-spatial-ordered},
\begin{align}
 [\rho_R(x),\rho_R(y)]
 &=\frac1{L^2}\sum_{m=1}^{M}m
   \bigl(\chi_m(x-y)-\chi_{-m}(x-y)\bigr)I,
 \label{eq:ch19-spatial-commutator}\\
 D_P(x,y)-D_P(y,x)
 &=\frac{2i}{L^2}\sum_{m=1}^{M}m
   \sin\!\left(\frac{2\pi m(x_0-y_0)}L\right)
 \label{eq:ch19-spatial-commutator-sine}
\end{align}
for the canonical root in the second line.  In particular, this
commutator expectation is independent of $c,s$ and of the chosen
normalized state on the exact CCR algebra.
\end{proposition}

\begin{proof}
The equal-sign generator commutators vanish.  In the expansion of
the commutator of the two sums, only the mixed terms survive:
\begin{align*}
 L^2[\rho_R(x),\rho_R(y)]
 &=\sum_{m,n=1}^{M}\bigl[
  \chi_{-m}(x)\chi_n(y)[C_{Rm},A_{Rn}]\\
 &\hspace{49mm}+
  \chi_m(x)\chi_{-n}(y)[A_{Rm},C_{Rn}]\bigr].
\end{align*}
By \eqref{eq:ch19-spatial-ccr},
$[C_{Rm},A_{Rn}]=-m\delta_{mn}I$ and
$[A_{Rm},C_{Rn}]=m\delta_{mn}I$.  Collapsing the deltas and using
\eqref{eq:ch19-spatial-character-laws} gives the first equation.
Applying $\omega_P(I)=1$ yields the second after
$e^{i\theta}-e^{-i\theta}=2i\sin\theta$.
Alternatively, subtracting \eqref{eq:ch19-spatial-ordered} with
$x,y$ interchanged produces the same expression because
$c^2-s^2=1$.  The direct CCR calculation demonstrates the stronger
fact: the commutator is already a scalar multiple of the unit
\emph{before} a state is chosen.
\end{proof}

\begin{proposition}[Symmetrized density correlation]
\label{prop:ch19-spatial-symmetrized}
Let
$S_P(x,y):=\tfrac12\omega_P(\{\rho_R(x),\rho_R(y)\})$.
Then for the canonical root
\begin{equation}
 \boxed{\displaystyle
 S_P(x,y)=\frac{c^2+s^2}{L^2}
 \sum_{m=1}^{M}m\cos\!\left(\frac{2\pi m(x_0-y_0)}L\right).}
 \label{eq:ch19-spatial-symmetrized}
\end{equation}
Set $g=(c-s)^2>0$, so $g^{-1}=(c+s)^2$.  Equivalently,
\begin{equation}
 c^2+s^2=\frac{g+g^{-1}}2,
 \qquad
 S_P(x,y)=\frac{g+g^{-1}}{2L^2}
 \sum_{m=1}^{M}m\cos\!\left(\frac{2\pi m(x_0-y_0)}L\right).
 \label{eq:ch19-spatial-symmetrized-g}
\end{equation}
\end{proposition}

\begin{proof}
Write $d=x_0-y_0$ and $\theta_m=2\pi md/L$.
The ordered formula takes the form
\begin{align*}
 D_P(x,y)
 &=\frac1{L^2}\sum_{m=1}^{M}m
   \left[s^2(\cos\theta_m-i\sin\theta_m)
        +c^2(\cos\theta_m+i\sin\theta_m)\right]\\
 &=\frac1{L^2}\sum_{m=1}^{M}m
   \left[(c^2+s^2)\cos\theta_m+i(c^2-s^2)\sin\theta_m\right]\\
 &=S_P(x,y)+\frac{i}{L^2}\sum_{m=1}^{M}m\sin\theta_m.
\end{align*}
The last term changes sign when $x$ and $y$ are exchanged,
whereas the first does not.  Their half-sum proves
\eqref{eq:ch19-spatial-symmetrized}, and the half-difference
agrees with Proposition~\ref{prop:ch19-spatial-commutator}.
Finally, $(c-s)(c+s)=1$ gives
$g^{-1}=(c+s)^2$ and hence
$g+g^{-1}=2(c^2+s^2)$.
\end{proof}

The decomposition is particularly transparent:
\begin{equation}
 \boxed{\displaystyle
 D_P(x,y)=S_P(x,y)
  +\frac{i}{L^2}\sum_{m=1}^{M}m
    \sin\!\left(\frac{2\pi m(x_0-y_0)}L\right).}
 \label{eq:ch19-spatial-ordered-realimag}
\end{equation}
The imaginary, antisymmetric contribution is fixed by the CCR; the
interaction-dependent parameter $g$ affects the real symmetric part.
This is an \emph{exact finite-mode} separation, not an assertion that
the imaginary part becomes negligible at large distance or large $L$.

\subsection{Coincident points, positivity, and what positivity means}

\begin{theorem}[Diagonal evaluation and positivity]
\label{thm:ch19-spatial-diagonal}
Under the hypotheses of Theorem~\ref{thm:ch19-spatial-ordered},
for every $x\in\Lambda_L$,
\begin{equation}
 \boxed{\displaystyle
 D_P(x,x)=\frac{c^2+s^2}{L^2}\frac{M(M+1)}2
           =\frac{g+g^{-1}}{2L^2}\frac{M(M+1)}2.}
 \label{eq:ch19-spatial-diagonal}
\end{equation}
This value is real and nonnegative.  It is strictly positive when
$M\geq1$ (and $h>0$, as assumed), whereas $M=0$ gives zero.
\end{theorem}

\begin{proof}
At $x=y$, every character in \eqref{eq:ch19-spatial-ordered}
is $\chi_{\pm m}(0)=1$, so
\[
 D_P(x,x)=\frac{c^2+s^2}{L^2}\sum_{m=1}^{M}m.
\]
The elementary identity $\sum_{m=1}^{M}m=M(M+1)/2$ proves the
formula, including the empty-sum case.  Since $c^2-s^2=1$,
\[
 c^2+s^2=1+2s^2\geq1.
\]
The denominator $L^2$ is positive for $h>0$.  Therefore the right
side is strictly positive precisely when the cutoff contains at least
one positive mode.  No restriction such as $M<h$ is used.
\end{proof}

This diagonal is the second moment of the Hermitian oscillatory
density in the abstract state:
$D_P(x,x)=\omega_P(\rho_R(x)^*\rho_R(x))$.
Its finiteness results from the \emph{mode cutoff} $M$, not from a
finite-dimensional oscillator representation.  Increasing $M$ at
fixed $L$ makes the diagonal grow quadratically; thus its finite
value does not, by itself, define a cutoff-independent continuum
coincident-point observable.

There is a stronger, more useful positivity statement than
$D_P(x,x)\geq0$.  A complex two-point function is a
\emph{positive-semidefinite kernel} if for any finite collection of
sites $x_1,\ldots,x_N$ and complex numbers $z_1,\ldots,z_N$,
\begin{equation}
 \sum_{a,b=1}^{N}\overline{z_a}\,D_P(x_a,x_b)\,z_b\geq0.
 \label{eq:ch19-spatial-psd}
\end{equation}
This inequality is an inequality of \emph{real} numbers: its left
side is automatically real.

\begin{proposition}[Positive-semidefinite ordered kernel]
\label{prop:ch19-spatial-psd}
The kernel $D_P$ of Theorem~\ref{thm:ch19-spatial-ordered} satisfies
\eqref{eq:ch19-spatial-psd}; it is also Hermitian,
$D_P(y,x)=\overline{D_P(x,y)}$.
\end{proposition}

\begin{proof}
For the chosen sites and coefficients, set
$B=\sum_{b=1}^{N}z_b\rho_R(x_b)$.
Each density is self-adjoint, so
$B^*=\sum_{a=1}^{N}\overline{z_a}\rho_R(x_a)$.
Consequently,
\begin{align*}
 \sum_{a,b}\overline{z_a}D_P(x_a,x_b)z_b
 &=\omega_P\!\left(
    \sum_{a,b}\overline{z_a}z_b\rho_R(x_a)\rho_R(x_b)
    \right)\\
 &=\omega_P(B^*B)\geq0,
\end{align*}
by positivity of the abstract functional.  Hermiticity follows
from either $\omega_P(B^*)=\overline{\omega_P(B)}$ or the explicit
formula \eqref{eq:ch19-spatial-ordered} and
$\chi_{-m}(d)=\overline{\chi_m(d)}$.
\end{proof}

Equivalently, \eqref{eq:ch19-spatial-ordered} expresses the matrix
$D_P(x_a,x_b)$ as a sum of rank-one Hermitian matrices with
nonnegative weights $ms^2/L^2$ and $mc^2/L^2$.  This gives a second
way to see positivity and shows its rank is at most $2M$, possibly
less because lattice Fourier characters can coincide.  Positive
semidefiniteness does \emph{not} mean that an off-diagonal entry is
real or positive: in general $D_P(x,y)$ is complex, and even its
symmetrized value $S_P(x,y)$ may be negative.

\subsection{Worked finite example: eight sites and two positive modes}

Take an eight-site periodic lattice and two abstract positive modes,
\begin{equation}
 L=8,\quad h=4,\quad M=2,\quad
 \zeta=e^{i\pi/4},\qquad c=\frac53,\quad s=\frac43.
 \label{eq:ch19-spatial-example-data}
\end{equation}
These choices satisfy $c^2-s^2=1$ and
$g=(c-s)^2=1/9$.  The two retained frequencies are distinct and
below the Nyquist mode $m=h=4$.  The normalized density at $x=0$
is explicitly
\[
 \rho_R(0)=\frac18(C_{R1}+A_{R1}+C_{R2}+A_{R2}),
\]
while at $x=1$ the corresponding four generators are weighted by
$\zeta^{-1},\zeta,\zeta^{-2},\zeta^2$, respectively.  Only the
pairings of opposite adjoint type and equal mode index survive when
the two densities are multiplied and evaluated.

For the separation $d=x_0-y_0$, the ordered kernel becomes
\begin{align}
 D_P(d)
 &=\frac1{64}\sum_{m=1}^{2}m
   \left(\frac{16}{9}e^{-i m\pi d/4}
        +\frac{25}{9}e^{i m\pi d/4}\right)\notag\\
 &=\frac1{576}\left[
  41\left(\cos\frac{\pi d}{4}
            +2\cos\frac{\pi d}{2}\right)
  +9i\left(\sin\frac{\pi d}{4}
            +2\sin\frac{\pi d}{2}\right)\right].
 \label{eq:ch19-spatial-example-finite}
\end{align}
In particular,
\begin{align}
 D_P(0)&=\frac{41}{192},
 \label{eq:ch19-spatial-example-diagonal}\\
 D_P(1)&=\frac{41\sqrt2}{1152}
          +i\,\frac{\sqrt2+4}{128},
 \label{eq:ch19-spatial-example-d1}\\
 D_P(-1)&=\frac{41\sqrt2}{1152}
          -i\,\frac{\sqrt2+4}{128}.
 \label{eq:ch19-spatial-example-dminus1}
\end{align}
Thus the finite example checks all three operations independently:
\begin{align*}
 S_P(1)&=\frac{D_P(1)+D_P(-1)}2
       =\frac{41\sqrt2}{1152},\\
 \omega_P([\rho_R(1),\rho_R(0)])
 &=D_P(1)-D_P(-1)
   =i\,\frac{\sqrt2+4}{64}.
\end{align*}
The commutator is independent of $c,s$: evaluating the same expression
with the free pair $(c,s)=(1,0)$ gives the \emph{same} imaginary
difference, although its real symmetric part changes.
For example, in that free pair,
\begin{equation}
 D_0(1)=\frac{\sqrt2}{128}
          +i\,\frac{\sqrt2+4}{128},
 \qquad D_0(0)=\frac3{64}.
 \label{eq:ch19-spatial-example-free}
\end{equation}
Both states obey the same exact CCR; only their fluctuation
covariances differ.

For completeness, the $2\times2$ matrix associated with the sites
$0$ and $1$ is
\[
 \begin{pmatrix}
  41/192&D_P(-1)\\
  D_P(1)&41/192
 \end{pmatrix}.
\]
It is Hermitian and positive semidefinite by
Proposition~\ref{prop:ch19-spatial-psd}, even though the entries
$D_P(\pm1)$ are not real.  This is the correct interpretation of
positive spatial correlations in a noncommutative theory.

\subsection{What has, and has not, been established}

The passage from Section~\ref{sec:ch19-mode-correlators} to this
section is an exact finite Fourier reconstruction within the
\emph{abstract, uncompressed} CCR algebra.  The main formula is
valid for every integer $M\geq0$ and every pair of lattice sites,
under the stated root-of-unity and hyperbolic-pair conditions.
The commutator kernel depends only on the CCR and the chosen Fourier
normalization.  The symmetric kernel depends on the squeezing
parameter through $c^2+s^2=(g+g^{-1})/2$.  The diagonal is finite at
each cutoff and positive when $M\geq1$, and the full ordered kernel
is positive semidefinite in the sense of quadratic forms.

Three distinct extrapolations require additional work.  First,
a finite periodic character sum is not a proof of an infinite-volume
power law: any such statement needs a specified limit of $L$, $M$,
separations, and operator normalization.  Second, because the
oscillator carrier is uncompressed, the kernel is not automatically
a correlator of compressed finite-CAR density matrices; the next
comparison section must explicitly track maps and boundary
remainders.  Third, the oscillatory density used here omits a charge
zero mode, so the result should not be presented as the complete
microscopic density--density function without supplying that missing
sector and its cross terms.  These are limitations of the identified
observable, not failures of the finite algebraic calculation.

The corresponding statements in the supplied
\texttt{BosonizeStubs/Ch19Correlations.lean} snapshot are:
\begin{itemize}
 \item \texttt{ordered\_density\_kernel} and
       \texttt{ordered\_density\_kernel\_lattice};
 \item \texttt{density\_commutator\_expectation};
 \item \texttt{symmetrized\_density\_kernel};
 \item \texttt{density\_kernel\_diagonal\_pos}.
\end{itemize}
Their Phase~B status is not Core promotion, and standalone LaTeX
compilation does not replace independent mathematical review or
whole-book integration.
